\section{Thiết kế các thuật toán và mô-đun xử lý cốt lõi}
\label{sec:core_algorithms}

Trong hệ thống hỗ trợ ra quyết định lâm sàng và trợ lý y tế, việc đảm bảo tính chính xác khoa học, tính tất định và kiểm soát chặt chẽ ranh giới câu trả lời của mô hình ngôn ngữ lớn là yêu cầu tiên quyết. Phần này trình bày chi tiết các mô hình toán học, cấu trúc dữ liệu và thuật toán lõi được hiện thực hóa trong hệ thống.

% ==============================================================================
\subsection{Kiến trúc 7 Tầng Kiểm Soát Ranh Giới Hội Thoại và Hạn Chế Ảo Giác (Anti-Drift \& Anti-Hallucination Framework)}
\label{subsec:anti_hallucination_framework}

Để giải quyết vấn đề mô hình ngôn ngữ lớn tự suy diễn số liệu liều lượng hoặc bị lệch hướng khỏi ngữ cảnh điều trị buồng bệnh, hệ thống áp dụng mô hình kiểm soát 7 tầng liên hoàn:

\begin{enumerate}[leftmargin=2em]
    \item \textbf{Tầng 1: Bộ định tuyến ý định NLU và Ranh giới ngữ cảnh đóng (\textit{NLU Intent Router \& Bounded Clinical Context}):}
    Khi bác sĩ gửi câu hỏi, module \code{GeminiNLURouter} phân tích cú pháp và phân loại ý định vào 4 miền nghiệp vụ hợp lệ:
    \begin{equation}
        \text{Intent}(Q) \in \{\text{DRUG\_INQUIRY}, \text{DOSAGE\_ADJUSTMENT}, \text{INTERACTION\_CHECK}, \text{REGULATORY\_BHYT}\}
    \end{equation}
    Nếu câu hỏi chứa từ khóa phi lâm sàng, cố ý vượt quyền (Jailbreak/Prompt Injection) hoặc hỏi về các chủ đề ngoài phạm vi điều trị nội trú, hệ thống lập tức kích hoạt bộ lọc bảo vệ, từ chối trả lời và hướng người dùng quay lại hồ sơ bệnh án của ca bệnh.

    \item \textbf{Tầng 2: Chuẩn hóa thực thể hoạt chất quốc tế INN (\textit{Active Ingredient Normalizer}):}
    Để tránh việc tìm kiếm sai lệch do khác biệt thương mại giữa các nhà sản xuất, module \code{active_ingredient_normalizer.py} chuẩn hóa mọi danh từ thuốc:
    \begin{itemize}[leftmargin=1.5em]
        \item Khử bỏ hàm lượng và dạng bào chế bằng biểu thức chính quy (ví dụ: \code{Augmentin 1g/200mg viên nén} $\to$ \code{Augmentin}).
        \item Khử gốc muối hóa học (\textit{Salt Stripping}): Chuyển các gốc muối \textit{besylate}, \textit{tartrate}, \textit{hydrochloride} về dạng base tinh khiết.
        \item Ánh xạ biệt dược sang hoạt chất INN đối chiếu với 54.885 số đăng ký thuốc của Cục Quản lý Dược (ví dụ: \textit{Nexium} $\to$ \code{Esomeprazole}, \textit{Plavix} $\to$ \code{Clopidogrel}).
        \item Phân rã dạng thuốc phối hợp (\textit{Amoxicillin/Acid clavulanic} thành hai hoạt chất độc lập để kiểm tra tương tác chéo).
    \end{itemize}

    \item \textbf{Tầng 3: Corrective RAG (CRAG) và Truy xuất lai ghép (\textit{Hybrid Retrieval FTS5 BM25 + Dense Vector}):}
    Hệ thống kết hợp đồng thời tìm kiếm ngữ nghĩa sâu (Dense Retrieval) và tìm kiếm từ khóa chính xác (Lexical Retrieval) trên bộ nhớ RAM:
    \begin{itemize}[leftmargin=1.5em]
        \item \textbf{Điểm BM25 FTS5:} Đo lường tần suất xuất hiện chính xác của thuật ngữ y khoa:
        \begin{equation}
            \text{Score}_{\text{BM25}}(D, Q) = \sum_{i=1}^n \text{IDF}(q_i) \cdot \frac{f(q_i, D) \cdot (k_1 + 1)}{f(q_i, D) + k_1 \cdot \left(1 - b + b \cdot \frac{|D|}{\text{avgdl}}\right)}
        \end{equation}
        với các tham số thực nghiệm: $k_1 = 1.2$, $b = 0.75$.
        \item \textbf{Điểm Cosine Similarity:} Đo lường sự tương đồng ngữ nghĩa giữa vector nhúng của câu hỏi $\mathbf{q}$ và đoạn trích văn bản Dược thư $\mathbf{d}_i$:
        \begin{equation}
            S_{\text{vector}}(\mathbf{q}, \mathbf{d}_i) = \frac{\mathbf{q} \cdot \mathbf{d}_i}{\|\mathbf{q}\|_2 \|\mathbf{d}_i\|_2}
        \end{equation}
        \item \textbf{Điểm số tổng hợp:} Kết hợp theo tỷ lệ trọng số:
        \begin{equation}
            S_{\text{final}} = 0.65 \cdot S_{\text{vector}} + 0.35 \cdot S_{\text{BM25}}
        \end{equation}
        \item \textbf{Ngưỡng lọc CRAG (\textit{Confidence Threshold}):} Nếu $\max_i S_{\text{final}} < \theta_{\text{min}} = 0.72$, hệ thống xác định văn bản truy xuất chưa đủ độ tương đồng tin cậy; thay vì để LLM tự suy diễn, hệ thống chủ động thông báo "Chưa đủ cơ sở dữ liệu pháp quy" và dẫn chiếu bác sĩ mở trực tiếp bản PDF gốc.
    \end{itemize}

    \item \textbf{Tầng 4: Ràng buộc Prompting và Khóa nhiệt độ \texorpdfstring{($T = 0.0$)}{(T = 0.0)}:}
    Nhiệt độ sinh ngẫu nhiên của mô hình được thiết lập cố định bằng $0.0$ ($T = 0.0$, top-p = 1.0) nhằm hạn chế tính ngẫu nhiên, bảo đảm tính tất định và khả năng tái lập kết quả. System Prompt chứa điều khoản cấm ngoại suy (\textit{Strict No-Extrapolation Constraint}): mọi kết luận phải có căn cứ trong đoạn Context được cung cấp.

    \item \textbf{Tầng 5: Định vị nguồn và Trích dẫn số trang chính xác (\textit{Page Grounding with Direct Citations}):}
    Mọi thông tin liên quan đến liều dùng hoặc khuyến cáo bắt buộc phải đính kèm con trỏ tài liệu nguồn có thể kiểm chứng được, theo định dạng chuẩn:
    \begin{center}
        \code{[DTQG: Trang xxx]} \quad hoặc \quad \code{[TT20: Mục yyy]} \quad hoặc \quad \code{[QĐ5948: Cặp zzz]}
    \end{center}
    Giao diện người dùng cho phép bác sĩ nhấn vào con trỏ này để mở trực tiếp trang PDF tương ứng trong vòng $0.1\,\text{s}$.

    \item \textbf{Tầng 6: Bộ hậu kiểm đối chiếu số liệu định lượng (\textit{Self-RAG Numerical Verifier}):}
    Trước khi câu trả lời được phát luồng (Stream) đến giao diện bác sĩ, module \code{self_rag_verifier.py} trích xuất toàn bộ các số liệu định lượng:
    \begin{equation}
        \mathcal{N}_{\text{gen}} = \text{ExtractNumbers}(\text{Response})
    \end{equation}
    (bao gồm các giá trị liều mg, g, mcg; khoảng cách đưa thuốc theo giờ; tỷ lệ \%; và chỉ số eGFR). Đồng thời trích xuất tập số liệu xuất hiện trong đoạn tài liệu gốc:
    \begin{equation}
        \mathcal{N}_{\text{source}} = \text{ExtractNumbers}(\text{Context})
    \end{equation}
    Hệ thống kiểm tra điều kiện bao hàm nghiêm ngặt:
    \begin{equation}
        \mathcal{N}_{\text{gen}} \subseteq \mathcal{N}_{\text{source}}
    \end{equation}
    Nếu tồn tại bất kỳ số liệu nào $\exists x \in \mathcal{N}_{\text{gen}} \setminus \mathcal{N}_{\text{source}}$ (số liệu xuất hiện trong câu trả lời nhưng không có trong Dược thư gốc), hệ thống xác định đây là ảo giác nguy hiểm.

    \item \textbf{Tầng 7: Cắt mạch khẩn cấp (\textit{Circuit Breaker \& Safe Fallback}):}
    Khi điều kiện kiểm tra ở Tầng 6 bị vi phạm, luồng phát câu trả lời lập tức bị hủy bỏ (Abort SSE stream). Hệ thống thay thế bằng cảnh báo an toàn màu đỏ: \textit{"Phát hiện mâu thuẫn số liệu định lượng giữa mô hình AI và văn bản Dược thư. Yêu cầu bác sĩ tra cứu trực tiếp tài liệu gốc đính kèm."}
\end{enumerate}

% ==============================================================================
\subsection{Cổng kiểm duyệt an toàn xác định CDSS Sentinel \texorpdfstring{($O(1)$)}{(O(1))}}
\label{subsec:cdss_sentinel}

Module \code{cdss_sentinel.py} hoạt động độc lập với LLM, thực hiện kiểm tra tương tác thuốc chống chỉ định (theo danh mục Quyết định 5948/QĐ-BYT) theo phương pháp so khớp bảng băm xác định:

\begin{itemize}[leftmargin=1.5em]
    \item \textbf{Cấu trúc khóa băm đối xứng:} Để tối ưu hóa việc tra cứu hai chiều giữa hoạt chất $A$ và hoạt chất $B$, khóa tra cứu trong bảng băm lưu trữ trong RAM được chuẩn hóa theo thứ tự từ điển:
    \begin{equation}
        \text{Key}(A, B) = \min(\text{INN}_A, \text{INN}_B) + \text{"\_"} + \max(\text{INN}_A, \text{INN}_B)
    \end{equation}
    \item \textbf{Độ phức tạp tính toán:} Với đơn thuốc gồm $N$ hoạt chất, số lượng cặp tương tác cần kiểm tra là tổ hợp chập 2:
    \begin{equation}
        K = \binom{N}{2} = \frac{N(N-1)}{2}
    \end{equation}
    Do mỗi phép tra cứu trong bảng băm diễn ra trong thời gian $O(1)$, tổng thời gian kiểm tra của toàn bộ đơn thuốc là $O(N^2)$. Trong thực tế điều trị nội trú, một đơn thuốc hiếm khi vượt quá 15 thuốc ($N \le 15 \implies K \le 105$ cặp), tổng thời gian thực thi trên CPU luôn dưới $0.5\,\text{ms}$, hoàn toàn không gây trễ giao diện.
    \item \textbf{Cơ chế khóa chặn (Blocking Gate):} Khi phát hiện cặp tương tác thuộc danh mục Cấm (Contraindicated), API trả về mã lỗi HTTP 409 Conflict, lập tức vô hiệu hóa nút "Lưu đơn thuốc" và hiển thị thẻ cảnh báo cơ chế dược lý cùng hậu quả lâm sàng (xuất huyết não, kéo dài khoảng QT, suy tim cấp).
\end{itemize}

% ==============================================================================
\subsection{Thuật toán Cockcroft-Gault và Phân tầng Chỉnh liều Suy Thận KDIGO}
\label{subsec:cockcroft_gault}

Để phòng ngừa ngộ độc thuốc do suy giảm chức năng bài xuất qua thận, hệ thống tích hợp công thức tính độ thanh thải Creatinine ước tính (\textit{Creatinine Clearance} --- CrCl) kinh điển của Cockcroft và Gault (1976) \cite{cockcroft_gault_1976}:

\begin{equation}
    \text{CrCl} = \frac{(140 - \text{Tuổi}) \times \text{Cân nặng (kg)}}{72 \times S_{\text{Cr}} (\text{mg/dL})} \times \alpha
\end{equation}
trong đó:
\begin{itemize}[leftmargin=1.5em]
    \item $\text{Tuổi}$ tính theo năm; $\text{Cân nặng}$ tính theo kilôgam (kg).
    \item $S_{\text{Cr}}$ là nồng độ Creatinine huyết thanh đo được của người bệnh. Nếu dữ liệu xét nghiệm cung cấp theo đơn vị SI ($\mu\text{mol/L}$), hệ thống tự động quy đổi:
    \begin{equation}
        S_{\text{Cr}} (\text{mg/dL}) = \frac{S_{\text{Cr}} (\mu\text{mol/L})}{88.4}
    \end{equation}
    \item Hệ số hiệu chỉnh theo giới tính $\alpha$:
    \begin{equation}
        \alpha = \begin{cases} 
        1.00 & \text{nếu người bệnh là Nam} \\
        0.85 & \text{nếu người bệnh là Nữ (do khối lượng cơ thấp hơn)}
        \end{cases}
    \end{equation}
\end{itemize}

Hệ thống phân tầng chức năng thận theo 5 mức chuẩn quốc tế KDIGO 2024:
\begin{itemize}[leftmargin=1.5em]
    \item \textbf{G1} ($\text{CrCl} \ge 90\,\text{mL/min}$): Chức năng thận bình thường.
    \item \textbf{G2} ($\text{CrCl} = 60 - 89\,\text{mL/min}$): Suy thận nhẹ; kiểm tra cảnh báo thận trọng.
    \item \textbf{G3a--G3b} ($\text{CrCl} = 30 - 59\,\text{mL/min}$): Suy thận vừa; kích hoạt giảm liều hoặc giãn khoảng cách đưa thuốc (ví dụ Enoxaparin, Ciprofloxacin).
    \item \textbf{G4} ($\text{CrCl} = 15 - 29\,\text{mL/min}$): Suy thận nặng; khóa lựa chọn thuốc chống chỉ định (ví dụ Metformin do nguy cơ nhiễm toan lactic, Spironolactone).
    \item \textbf{G5} ($\text{CrCl} < 15\,\text{mL/min}$): Suy thận giai đoạn cuối/chạy thận nhân tạo; bắt buộc hội chẩn chuyên khoa thận.
\end{itemize}

% ==============================================================================
\subsection{Thuật toán Kiểm soát Trần Liều Tích Lũy 24 Giờ (Cumulative 24h Ceiling)}
\label{subsec:cumulative_dose_ceiling}

Sai sót quá liều do kê đơn nhiều chế phẩm có cùng hoạt chất là biến cố phổ biến khi điều trị đa khoa. Module \code{dose_ceiling_validator.py} tính tổng liều tích lũy của từng hoạt chất trong một ngày:

\begin{equation}
    D_{\text{total}}(A) = \sum_{k=1}^m \text{Dose}_k(A) \times \text{Frequency}_k
\end{equation}
trong đó $\text{Dose}_k(A)$ là hàm lượng của hoạt chất $A$ trong lần dùng thứ $k$, và $\text{Frequency}_k$ là số lần đưa thuốc trong 24 giờ.

Điều kiện kiểm định an toàn:
\begin{equation}
    D_{\text{total}}(A) \le D_{\text{max\_24h}}(A)
\end{equation}
Ví dụ với Paracetamol (Acetaminophen): $D_{\text{max\_24h}} = 4.000\,\text{mg/ngày}$ đối với người lớn có cân nặng $\ge 50\,\text{kg}$ và chức năng gan bình thường. Nếu người bệnh được kê đồng thời \code{Paracetamol 500mg} (uống 4 lần/ngày = $2.000\,\text{mg}$) và dịch truyền \code{Perfalgan 1g} (truyền 3 chai/ngày = $3.000\,\text{mg}$), tổng liều tích lũy là $5.000\,\text{mg/ngày} > 4.000\,\text{mg/ngày}$. Hệ thống lập tức phát báo động đỏ ngăn chặn ngộ độc hoại tử tế bào gan cấp tính.

% ==============================================================================
\subsection{Thuật toán Kiểm Định Xuất Toán BHYT Phân Tầng Theo Hạng Bệnh Viện}
\label{subsec:bhyt_audit_algorithm}

Căn cứ Thông tư 20/2022/TT-BYT, mỗi hoạt chất $A$ được phân định phạm vi thanh toán theo hạng bệnh viện cao nhất được phép chỉ định:

\begin{equation}
    \text{TierLevel} \in \{\text{Đặc biệt} (1), \text{Hạng I} (2), \text{Hạng II} (3), \text{Hạng III} (4), \text{Hạng IV} (5)\}
\end{equation}

Hệ thống so sánh hạng kỹ thuật của cơ sở khám chữa bệnh hiện tại ($H_{\text{current}}$) với hạng cho phép của thuốc ($H_{\text{allowed}}(A)$):
\begin{equation}
    \text{Status}(A) = \begin{cases}
    \text{ALLOWED (Được thanh toán)} & \text{nếu } H_{\text{current}} \le H_{\text{allowed}}(A) \\
    \text{REJECTION\_RISK (Nguy cơ xuất toán)} & \text{nếu } H_{\text{current}} > H_{\text{allowed}}(A)
    \end{cases}
\end{equation}
Khi phát hiện thuốc vượt tuyến, hệ thống không chỉ cảnh báo mà còn kích hoạt thuật toán gợi ý thay thế 1-Click: tìm kiếm trong đồ thị hoạt chất các thuốc cùng phân nhóm giải phẫu--điều trị ATC cấp 4--5 có $H_{\text{allowed}} \ge H_{\text{current}}$ và có sẵn trong danh mục tủ trực buồng bệnh.

% ==============================================================================
\subsection{Thuật toán Phân Bổ Xuất Kho Dược Theo Nguyên Tắc FEFO (First Expired, First Out)}
\label{subsec:fefo_algorithm}

Để giảm thiểu tổn thất tài chính do thuốc hết hạn trong tủ trực, module quản lý kho dược áp dụng thuật toán FEFO:

\begin{enumerate}[leftmargin=1.5em]
    \item Tập hợp tất cả các lô thuốc hiện có trong kho đối với hoạt chất yêu cầu:
    \begin{equation}
        \mathcal{B} = \{b_1, b_2, \dots, b_k\}
    \end{equation}
    trong đó mỗi lô $b_i$ có số lượng tồn kho $S(b_i)$ và hạn dùng $\text{ExpDate}(b_i)$.
    \item Sắp xếp tập $\mathcal{B}$ theo thứ tự tăng dần của hạn sử dụng:
    \begin{equation}
        \text{ExpDate}(b_1) \le \text{ExpDate}(b_2) \le \dots \le \text{ExpDate}(b_k)
    \end{equation}
    \item Khi nhận yêu cầu cấp phát số lượng $Q_{\text{req}}$, thuật toán lần lượt trừ tồn kho từ lô $b_1$:
    \begin{equation}
        q_i = \min(Q_{\text{req}} - \sum_{j=1}^{i-1} q_j, S(b_i))
    \end{equation}
    cho đến khi $\sum_{i} q_i = Q_{\text{req}}$.
\end{enumerate}

% ==============================================================================
\subsection{Thuật toán Băm Mật Khẩu Chống Tấn Công Brute-Force Argon2id}
\label{subsec:argon2id_security}

Để bảo vệ tài khoản nhân viên y tế và tuân thủ Luật An toàn thông tin mạng, hệ thống sử dụng thuật toán Argon2id (phiên bản lai ghép giữa Argon2d và Argon2i theo khuyến nghị của IETF RFC 9106 \cite{argon2id_rfc9106}):

\begin{equation}
    \text{Hash} = \text{Argon2id}(\text{Password}, \text{Salt}, t, m, p, L)
\end{equation}
với các tham số cấu hình:
\begin{itemize}[leftmargin=1.5em]
    \item Độ dài muối ngẫu nhiên ngầm định $\text{Salt}$: 16 bytes sinh bằng bộ sinh số ngẫu nhiên mật mã (\code{os.urandom}).
    \item Số vòng lặp thời gian (\textit{Time Cost}) $t = 3$ vòng.
    \item Bộ nhớ tiêu thụ (\textit{Memory Cost}) $m = 65.536\,\text{KiB}$ ($64\,\text{MB}$). Tham số bộ nhớ lớn giúp tăng độ phức tạp tính toán và nâng cao khả năng kháng cự trước các kỹ thuật tấn công dò vét bằng phần cứng chuyên dụng (GPU/ASIC).
    \item Mức độ song song hóa (\textit{Parallelism}) $p = 4$ luồng.
    \item Độ dài khóa băm đầu ra $L = 32$ bytes.
\end{itemize}

Kết hợp với cơ chế xác thực hai kênh (\textit{Dual-Channel Authentication}) truyền đồng thời HttpOnly Cookie chống tấn công XSS và Bearer Token Header chống mất phiên làm việc trên trình duyệt di động.
